\documentclass[../../../include/open-logic-section]{subfiles}

\begin{document}

\olfileid{sth}{story}{extensionality}
\olsection{Ekstensionalitas}

Prakara kapisan kang kudu diandharake yaiku himpunan dibedakake
dening !!{element}s-e. Luwih cetha:

\begin{axiom}[Ekstensionalitas]
Yen himpunan $A$ lan $B$ nduweni !!{element}s kang padha, mula $A$
lan $B$ iku himpunan kang padha.
\[
  \lforall[A][\lforall[B][(\lforall[x][(x \in A \liff x \in B)] \lif
  \eq[A][B])]]
\]
\end{axiom}

Kita wis nganggep iki lumaku ing saindhenging \olref[sfr][][]{part}.
Nanging prinsip iki perlu dibaleni. Aksioma Ekstensionalitas
ngandharake gagasan dhasar yen sawijining himpunan ditemtokake
dening !!{element}s-e. (Mula himpunan bisa dibedakake saka
\emph{konsep}, amarga barang-barang kang persis padha bisa
kacakup ing akeh konsep kang beda.)

Yagene prinsip iki ditampa? Lumrah yen kita nganggep saben
penolakan marang Ekstensionalitas minangka kaputusan kanggo
ninggalake apa wae kang isih bisa diarani \emph{teori himpunan}.
Teori himpunan pancen teori babagan koleksi ekstensional.

Nanging tantangan kang sejati ing \olref[sth][][]{part} yaiku
nemtokake prinsip kang ngandharake \emph{himpunan endi kang ana}.
Lan ternyata siji-sijine wangsulan kang katon temen-temen
``cetha'' tumrap pitakonan iki bisa dibuktekake luput.

\end{document}
